\ChapterPage{Páginas 5--6 \textbullet{} Transcripción}{Entrevista en la aduana}
\textsc{Registrada por el escribiente J. Vale, 4 de octubre de 1928. La
siguiente transcripción se reconstruye a partir de notas; las pausas
indican fallos de la aguja del gramófono.}

\medskip
\noindent\textbf{VALE:} Diga su nombre para el acta.\\
\textbf{BELL:} Mara Bell.\\
\textbf{VALE:} ¿Y el doctor?\\
\textbf{VOSS:} Elian Voss. Estoy aquí como testigo, no como médico.\\
\textbf{VALE:} ¿Entraron en el sótano?\\
\textbf{BELL:} Bajamos la escalera. El sótano no estaba al final.

\begin{quote}
\textbf{VALE:} Explique esa diferencia.\\
\textbf{BELL:} La escalera terminaba en una puerta. Daba a un pasadizo que
no cabía bajo el edificio.\\
\textbf{VOSS:} La arquitectura solo es imposible si se supone que el
edificio se apoya en terreno corriente.\\
\textbf{VALE:} ¿En qué podría apoyarse?\\
\textbf{VOSS:} En unos cimientos, quizá. O en el recuerdo de unos cimientos.
\end{quote}

\textbf{VALE:} ¿Vieron alguna criatura?\\
\textbf{BELL:} Vimos una sombra en el vidrio.\\
\textbf{VALE:} ¿Un reflejo?\\
\textbf{BELL:} Se movió antes que nosotros.\\
\textbf{VOSS:} Los paneles no reflejaban la sala. Reflejaban la marea a una
distancia que no podíamos medir.
